\scanpage{49}
est majoré par $(n-1)d$, (au lieu de $(n-1)m^2$) où $d=\dim G$. \deleted{\unsure{Conséq.}} D'ailleurs, si $v_1,v_2,\ldots v_n$ sont donnés, \deleted{\unsure{les}} \deleted{prenant} définissant $g_1,\ldots,g_n\in G(\widehat{\mathbf Z})$ alors prenant $\sigma_2\ldots\sigma_n\in\operatorname{Gal}(k/k_0)$ (où $k_0$ est un corps de définition \inserted{de type fini} de $M$) et $v'_1,v'_2,\ldots v'_n$ avec $v'_i=\sigma_i v_i$ $1\leq i\leq n$, on trouve les valeurs $\unsure{g_1},\bar\sigma_2g_2,\ldots\bar\sigma_ng_n$, où $\bar\sigma_i$ est l'image dans $G(\widehat{\mathbf Z})$. \deleted{Car}

\margin{N.B. On le prendra $G$ ``\uline{\unsure{géom.}}'' donc \uline{connexe}}
Admettant la \uuline{conjecture de Tate renforcée}, $\operatorname{Gal}(k/k_0)\mathbin{\text{---}}G(\widehat{\mathbf Z})$ \uline{a une image qui est d'indice fini}. Il s'ensuit que quitte à remplacer les $v_i$ par des $\sigma_i v_i$ convenables, on peut \deleted{\unsure{trouver}} \inserted{obtenir} $g_2\ldots g_n$ gén. indépendants pour \uline{$\forall$} $\ell$, si on admet ceci :
\[
\left[\begin{array}{l}
\text{Soit }G\text{ grp alg. affine connexe sur }\mathbf Q\\
U\text{ sous-groupe ouvert cpt de }G(\widehat{\mathbf Q})\\
P\inserted{=U.a}\text{ une \unsure{cl.} dans }G(\widehat{\mathbf Q})\text{ mod }U\\
\text{Alors il existe dans }P\text{ un élément }\xi\text{ dont toutes les}\\
\text{composantes sont génériques rel./}\mathbf Q
\end{array}\right.
\]
\mspage{48}
